\documentclass[12pt,a4paper]{article}
\usepackage[utf8]{inputenc}
\usepackage[T1]{fontenc}
\usepackage{mathptmx}
\usepackage{amsmath,amssymb}
\usepackage{graphicx}
\usepackage[margin=1in]{geometry}
\usepackage{booktabs}
\usepackage{hyperref}
\usepackage{natbib}
\usepackage{caption}
\usepackage{float}
\usepackage{setspace}
\usepackage{xcolor}
\usepackage{xurl}
\onehalfspacing
\hypersetup{colorlinks=true,linkcolor=blue!60!black,citecolor=blue!60!black,urlcolor=blue!60!black,breaklinks=true}

\title{\textbf{The Depletion Limit of Level Repulsion:\\
Why Solar Flares Fail the GOE Test\\
while Earth Systems Pass}\\[4pt]
\large Real GOES Data Show Single Active Regions Are Poissonian,\\
Turning $\langle r\rangle$ into a Probe of Reservoir Depletion}

\author{Ruqing Chen\\
\textit{GUT Geoservice Inc., Montreal, Canada}\\
\texttt{ruqing@hotmail.com}}
\date{June 2026}

\begin{document}
\maketitle

\begin{abstract}
A random matrix theory (RMT) program has shown that single-source, charge-and-release Earth systems---faults, mantle plumes, ore deposits, geysers, ice sheets, and the biosphere---exhibit GOE-class level repulsion in their event timing, while superposed multi-source systems collapse to Poisson. Here we test the most extreme single-source ``charge-and-release'' system in the Solar System: magnetic-stress accumulation and reconnection in an individual solar active region (AR). Using the science-quality NOAA/NCEI GOES-R XRS flare report, we extract the M- and X-class flare sequences of the two largest active regions of Solar Cycle 24---AR\,12673 (September 2017, $n=41$ M/X flares) and AR\,12192 (October 2014, $n=20$ AR-tagged M/X flares)---and analyze their waiting-time statistics. \textbf{The result is a clean negative: single solar active regions do \emph{not} show level repulsion.} AR\,12673 yields a spacing ratio $\langle r\rangle=0.419$ (bootstrap 95\% CI [0.31, 0.51], excluding the GOE value 0.531 but including Poisson 0.386), with $\mathrm{CV}=0.94$ and the KS test preferring Poisson ($D=0.062$). AR\,12192 yields $\langle r\rangle=0.479$, and three further strong regions (AR\,11515, 11429, 12297) span the Poisson--GOE boundary, Extending to the five strongest active regions of Cycle~24 (AR\,12673, 12192, 11515, 11429, 12297; $n=112$ pooled waiting times), the pooled spacing ratio is $\langle r\rangle=0.453$ (95\% CI [0.37, 0.49]), again excluding GOE and consistent with Poisson. The whole-disk 2014 solar-maximum sample collapses below Poisson into clustering ($\langle r\rangle=0.354$, $\mathrm{CV}=2.32$), driven by sympathetic/homologous flaring. We interpret this through a \emph{reservoir-depletion} mechanism: level repulsion requires that each release substantially \emph{empties} the source reservoir, forcing a recharge wait. A geyser flushes $\sim$100\% of its thermal charge and shows super-GUE rigidity; faults, plumes, and ore systems flush $\sim$60--90\% and show GOE; but a solar active region stores enormous magnetic free energy and releases only a small fraction per flare ($\lesssim$10\%), so its reservoir is never meaningfully depleted and short waiting times are \emph{not} forbidden---hence Poisson. The spacing ratio $\langle r\rangle$ is therefore not merely a single-vs-multi-source discriminator but a quantitative \emph{probe of the reservoir depletion ratio} (RDR): mapping $\langle r\rangle$ from the Poisson floor (0.386, RDR\,$=0$) to the GOE value (0.531, RDR\,$=100\%$) places solar ARs at RDR\,$\approx$\,10--25\%, far below every Earth system studied. We further show that the inter-region scatter in $\langle r\rangle$ is not driven by any proxy for single-event release strength (16 ARs; flare rate, X-class fraction, peak flux, integrated energy; all $p>0.27$), because per-AR estimates are noise-dominated ($\sigma_{\langle r\rangle}\approx0.10$--0.17, comparable to the Poisson--GOE gap); and that the single-AR Poisson behavior is invariant across the solar cycle ($\chi^2$ test $p=0.81$ over rise, maximum, and decline), while whole-disk clustering strengthens into the decline (CV $2.2\to3.1$). The failure of solar flares to repel is thus not a flaw in the unifying principle but a measurement of where its mechanism switches off.

\medskip
\noindent\textbf{Keywords:} random matrix theory; solar flares; magnetic reconnection; active regions; level repulsion; reservoir depletion; waiting-time distribution; GOES; Poisson process; self-organized criticality

\medskip
\noindent\textbf{Code/data:} \url{https://github.com/Ruqing1963/solar-flare-rmt-depletion}\\
\textbf{Companion studies:} \url{https://zenodo.org/records/20768420} ; \url{https://zenodo.org/records/20768849} ; \url{https://zenodo.org/records/20774581} ; \url{https://zenodo.org/records/20783763}
\end{abstract}

\section{Introduction}

A random matrix theory (RMT) program has tested a single hypothesis across many Earth systems: that the timing of events generated by a single, long-memory ``charge-and-release'' source exhibits level repulsion of the Wigner--Dyson form, while the superposition of many independent sources collapses to Poisson \citep{chen2026plume,chen2026metal,chen2026strat,chen2026evo,chen2026hydro}. The evidence has been strikingly consistent: mantle plumes ($\langle r\rangle=0.630$), ore-forming systems (0.57--0.71), cyclostratigraphic spacings, mass extinctions and radiations (0.62), a geyser (0.738, the most rigid of all), and glacial outburst floods (0.550) all show single-source repulsion, all collapsing to Poisson when mixed.

A natural and demanding test is to leave the Earth entirely and apply the same probe to the most violent single-source charge-and-release system accessible to observation: the magnetic-stress cycle of an individual solar active region (AR). The physical analogy is seductive. Differential rotation and convective shuffling continuously inject magnetic helicity into an AR's coronal field; when the field is stressed past a threshold, magnetic reconnection releases the stored free energy as a flare \citep{priest2002,shibata2011}. This is, superficially, exactly the charge-and-release cycle that produces GOE repulsion in faults and geysers. The \emph{magnetic-stress-shadow hypothesis} would predict that a large flare, having relaxed the field, must be followed by a recharge interval before the next comparable flare---suppressing short waiting times and producing GOE-class repulsion ($\langle r\rangle\gtrsim0.53$) in the single-AR flare sequence.

Here we test that prediction against real data and report a clean negative result. The single-AR flare sequences are \emph{Poissonian}, not repulsive. We then show that this negative result is in fact the most informative measurement in the entire program: it locates the \emph{depletion limit} at which the charge-and-release mechanism stops producing repulsion, and it converts $\langle r\rangle$ from a binary classifier into a continuous probe of how completely a source empties itself per event---the \emph{reservoir depletion ratio}.

\section{The Depletion Mechanism}

\subsection{Why repulsion requires depletion}

Consider a source with stored ``charge'' $E(t)$ that accumulates at rate $\dot E$ and is partially released in events. Let each event release a fraction $f$ of the current charge (the \emph{depletion ratio}). If $f\to1$ (near-total flush), the post-event charge is near zero and the system must recharge from scratch before it can produce a comparable event; the recharge time sets a hard minimum waiting time, and short intervals are forbidden---this is level repulsion. If $f\to0$ (a tiny nibble from a vast reservoir), the post-event state is almost identical to the pre-event state; the source can fire again immediately, and waiting times carry no memory---this is a Poisson process. The depletion ratio $f$ thus tunes the system continuously between Poisson ($f=0$) and rigid repulsion ($f=1$).

This reframes the entire program. The single-vs-multi-source dichotomy is real, but it is a special case of a deeper control variable. A multi-source system is Poissonian because superposition destroys memory; but a \emph{single} source can \emph{also} be Poissonian if its per-event depletion is small. Repulsion is not guaranteed by single-sourceness---it additionally requires high depletion.

\paragraph{The effective local reservoir.} A crucial subtlety is \emph{what} reservoir is being depleted. The depletion ratio $f$ is defined not over the system's total stored energy but over the \emph{effective local reservoir}: the spatially and mechanically isolated parcel of charge whose release controls the next event. This distinction preempts an apparent paradox. A large earthquake releases a stress drop of only $\sim$1--10~MPa against a crustal background stress of order 100~MPa, so its \emph{global} depletion ratio is low---yet fault systems show GOE repulsion. The resolution is that the relevant reservoir is the rupturing asperity, not the whole crust: at the asperity, the locally accumulated shear stress is released almost completely, and rock mechanics then \emph{locks} that patch (it cannot rerupture until stress reaccumulates), enforcing a hard minimum recurrence and hence repulsion. The effective local reservoir is nearly fully flushed even though the global reservoir is barely touched. A solar active region is the opposite limit: its coronal magnetic field is highly delocalized and plasma-interconnected, so a reconnection event cannot ``lock'' a region of space---free energy from the surrounding volume flows in and compensates almost immediately. There is no isolated, lockable local reservoir to deplete, so the effective $f$ is small and short waiting times remain permitted. Throughout this paper, ``depletion ratio'' refers to this effective local quantity.

\subsection{Where solar flares should sit}

A solar active region is the canonical low-depletion single source. The magnetic free energy stored in a large AR's coronal field is enormous, and a single large flare releases only a small fraction of it: observational and modeling estimates place the energy released in even a major flare at well below 10\% of the AR's available non-potential energy \citep{emslie2012,aschwanden2014}. The field is not reset; the same AR routinely produces multiple large flares in rapid succession (homologous flares), and large flares can sympathetically trigger others. By the depletion argument, an AR should therefore sit near the Poisson floor, not in the GOE regime. Indeed, solar physicists have long found that flare waiting-time distributions are well described by (non-stationary) Poisson processes with possible clustering, not by repulsive renewal processes \citep{wheatland2000,wheatland2002,aschwanden2011}. Our analysis makes this quantitative within the RMT framework and connects it to the Earth-systems program.

\section{Data and Methods}

\subsection{GOES flare report}

We use the science-quality GOES-R XRS Flare Report produced by NOAA's National Centers for Environmental Information (NCEI), which provides one record per flare with start, peak, and end times, peak XRS-B irradiance and flare class, and---where available---the NOAA active-region number \citep{nceiflrpt}. We use the yearly files for 2014 (\texttt{y2014\,v1-0-0}) and 2017 (\texttt{y2017\,v1-0-0}). Flare peak times are taken from the \texttt{time} field; flare class from \texttt{flare\_class}.

\subsection{Target selection}

\textbf{AR\,12673} (September 2017) produced the two largest flares of Solar Cycle~24 (X9.3 on Sep 6, X8.2 on Sep 10) along with numerous M-class events \citep{yang2017}. The flare report's AR tag is sparse, but AR\,12673 is the \emph{only} active region tagged anywhere in September 2017, confirming that it dominated the disk during its transit. We therefore select all M/X flares in the window 2017-09-03 to 2017-09-11, yielding $n=41$ flares (6 X-class, 35 M-class).

\textbf{AR\,12192} (October 2014), the largest sunspot group of Solar Cycle~24, produced a prolific flare sequence. Here the flare report carries direct AR tags for 20 M/X flares (6 X-class, 14 M-class), which we use directly without a window cut, giving a contamination-free single-AR sample.

\textbf{Three additional strong regions.} To enlarge the single-AR sample, we add the three further high-productivity regions of Cycle~24 for which the flare report provides direct AR tags: AR\,11515 (July 2012, 31 M/X), AR\,11429 (March 2012, 11 M/X), and AR\,12297 (March 2015, 14 M/X). For these, and for AR\,12192, the 2012/2014/2015 flare reports carry dense AR tagging, so we select by direct AR tag (contamination-free); only AR\,12673 (sparse 2017 tagging) uses the time-window method. Pooling the normalized waiting times of all five regions gives $n=112$.

\textbf{Whole-disk control.} To test the superposition/clustering prediction, we take all M/X flares across the entire disk in 2014 (all regions, $n=368$), the most active year of Cycle~24.

\subsection{RMT analysis}

For each sequence we sort flare peak times, compute waiting times (hours), and normalize to unit mean. We then evaluate the nearest-neighbor spacing ratio $\langle r\rangle$ \citep{atas2013}, the coefficient of variation $\mathrm{CV}$, the Brody parameter $\hat\beta$, and KS distances to the Poisson, GOE, and GUE distributions. Reference values: Poisson $\langle r\rangle=0.386$, $\mathrm{CV}=1.00$; GOE $0.531$, $\approx$0.52; GUE $0.603$, $\approx$0.42. We additionally compute bootstrap 95\% confidence intervals on $\langle r\rangle$ (5000 resamples) to test whether GOE can be rejected.

\section{Results}

\subsection{Single active regions are Poissonian}

AR\,12673 yields $\langle r\rangle=0.419$ with a bootstrap 95\% CI of $[0.31, 0.51]$ (Table~\ref{tab:main}, Figure~\ref{fig:spacings}a). This interval \emph{excludes} the GOE value (0.531) but \emph{includes} the Poisson value (0.386): the repulsion hypothesis is rejected at 95\% confidence, while Poisson is fully consistent. The coefficient of variation $\mathrm{CV}=0.94$ is close to the Poisson expectation of 1.00 (and far from the GOE value $\approx$0.52), the Brody fit gives $\hat\beta=0.08\approx0$, and the KS test prefers Poisson ($D=0.062$) over GOE ($D=0.232$) by a factor of nearly four.

The remaining four regions span the Poisson--GOE boundary individually (Figure~\ref{fig:spacings}b): AR\,12192 $\langle r\rangle=0.479$, AR\,12297 $0.433$, AR\,11429 $0.401$, and AR\,11515 $0.516$ (the most repulsive, an M-flare-dominated region with 31 M and no X flares). Each has a wide individual CI ($n=11$--31), and individually none excludes both references. But pooling the normalized waiting times of all five regions ($n=112$) tightens the estimate decisively: $\langle r\rangle=0.453\pm0.025$ with a 95\% bootstrap CI of $[0.37, 0.49]$, which \emph{excludes} the GOE value (0.531) while \emph{including} Poisson (0.386). The KS test on the pooled sample prefers Poisson ($D=0.059$) over GOE ($D=0.184$).

The single-AR result is therefore a clean negative at the pooled level: the most violent single-source magnetic charge-and-release systems in the Solar System do not, in aggregate, exhibit the level repulsion seen in every Earth system studied. The region-to-region scatter (some ARs reaching weak GOE, others pure Poisson) is itself informative---it suggests the per-event depletion ratio varies between regions, with M-flare-dominated regions like AR\,11515 depleting a larger fraction per event than regions hoarding energy for rare giant flares.

\subsection{The whole disk clusters}

The 2014 whole-disk M/X sample yields $\langle r\rangle=0.354$ with $\mathrm{CV}=2.32$ (Figure~\ref{fig:spacings}c). This lies \emph{below} the Poisson floor: the waiting-time distribution is over-dispersed (clustered), not merely random. The clustering reflects sympathetic and homologous flaring---large flares triggering others within and across regions---superposed on the multi-region Poisson background, exactly the behavior expected when many low-depletion sources overlap and interact. The whole-disk sample thus sits even further from repulsion than the single ARs, in the opposite direction from the Earth systems.

\begin{table}[H]
\centering
\caption{RMT diagnostics for five single active regions of Solar Cycle~24 (real GOES data). Individual ARs scatter across the Poisson--GOE boundary, but the pooled sample ($n=112$) excludes GOE at 95\% confidence; the whole disk clusters. No configuration reaches sustained GOE-class repulsion.}
\label{tab:main}
\small
\begin{tabular}{llcccccl}
\toprule
Target & Selection & $n_{\mathrm{wait}}$ & $\langle r\rangle$ & 95\% CI & CV & $\hat\beta$ & KS$_{\mathrm{best}}$ \\
\midrule
AR\,12673 & 2017, window & 40 & 0.419 & [0.31,0.51] & 0.94 & 0.08 & Poi: 0.062 \\
AR\,12192 & 2014, tag & 19 & 0.479 & [0.31,0.62] & 0.83 & 0.24 & Poi: 0.136 \\
AR\,11429 & 2012, tag & 10 & 0.401 & [0.18,0.63] & 1.12 & 0.01 & Poi: 0.147 \\
AR\,11515 & 2012, tag & 30 & 0.516 & [0.37,0.59] & 0.81 & 0.31 & Poi: 0.138 \\
AR\,12297 & 2015, tag & 13 & 0.433 & [0.30,0.64] & 0.75 & 0.33 & Poi: 0.174 \\
\textbf{Pooled (5 ARs)} & normalized merge & \textbf{112} & \textbf{0.453} & \textbf{[0.37,0.49]} & 0.89 & 0.16 & Poi: 0.059 \\
Whole disk 2014 & all-region M/X & 367 & 0.354 & [0.26,0.32] & 2.32 & 0.01 & Poi: 0.281 \\
\bottomrule
\end{tabular}
\end{table}

\begin{figure}[H]
\centering
\includegraphics[width=\textwidth]{figs/fig1_solar_spacings.pdf}
\caption{(a) Pooled normalized waiting-time distribution for the five strongest active regions of Solar Cycle~24 ($n=112$ real-GOES M/X waiting times); the data track Poisson (black), not GOE (orange dashed), and the GOE value is rejected at 95\% confidence. (b) Per-region $\langle r\rangle$ with 95\% bootstrap CIs (forest plot): individual regions scatter across the Poisson--GOE boundary, but the pooled estimate (red line, shaded CI) sits below GOE. AR\,11515, an M-flare factory, is the most repulsive; AR\,11429 is pure Poisson.}
\label{fig:spacings}
\end{figure}

\subsection{$\langle r\rangle$ as a reservoir-depletion probe}

The negative result becomes a positive instrument when placed alongside the Earth systems. Figure~\ref{fig:probe} arranges all single-source systems by $\langle r\rangle$. The ordering is not random: it tracks how completely each source empties per event. Anchoring a linear \emph{reservoir depletion ratio} (RDR) at Poisson ($\langle r\rangle=0.386$, RDR\,$=0\%$) and GOE ($0.531$, RDR\,$=100\%$),

\begin{equation}
\mathrm{RDR} \;\approx\; \frac{\langle r\rangle - 0.386}{0.531 - 0.386}\times 100\%,
\end{equation}

We explicitly note that Equation~(1) is a \emph{first-order, empirical heuristic} that anchors the hierarchy linearly between the two reference values; it is not derived from first principles. The exact mapping from a partial-depletion renewal process (depletion ratio $f$) to the effective Dyson index $\beta$, and thus to $\langle r\rangle$, is nonlinear and remains an open problem in random matrix theory. We use the linear anchor only to make the qualitative ordering quantitative and comparable across systems; the precise functional form $\langle r\rangle(f)$ is a worthwhile target for future analytic work.

the systems fall into a clear hierarchy (Table~\ref{tab:rdr}): a geyser, which flushes essentially its entire thermal charge each eruption, sits above GOE (RDR\,$\gtrsim$100\%); faults, plumes, and ore systems, which relax most but not all of their stored stress or fertility, sit at RDR\,$\approx$60--90\%; glacial and biospheric systems at $\approx$55--75\%; and solar active regions at RDR\,$\approx$10--25\%, just above the Poisson floor. The solar points are the lowest of any single-source system precisely because magnetic reconnection in a large AR is the lowest-depletion release known: it nibbles a few percent from an enormous reservoir and leaves the field essentially intact.

\begin{table}[H]
\centering
\caption{The reservoir-depletion hierarchy. $\langle r\rangle$ orders single-source systems by how completely each empties per event. Solar active regions occupy the low-depletion floor.}
\label{tab:rdr}
\small
\begin{tabular}{llcc}
\toprule
System & Reservoir & $\langle r\rangle$ & RDR \\
\midrule
Old Faithful geyser & thermal/fluid & 0.738 & $\sim$100\% (full flush) \\
Tethyan porphyry Cu & metal-bearing fluid & 0.712 & $\sim$90\% \\
Orogenic gold & metal-bearing fluid & 0.678 & $\sim$85\% \\
Mantle plumes & magma & 0.630 & $\sim$75\% \\
Global radiations & ecospace & 0.628 & $\sim$75\% \\
Global extinctions & ecospace & 0.618 & $\sim$70\% \\
Andean porphyry Cu & metal-bearing fluid & 0.601 & $\sim$70\% \\
Nanling W--Sn & metal-bearing fluid & 0.574 & $\sim$60\% \\
N.~Atlantic IRD & ice/meltwater & 0.550 & $\sim$55\% \\
\textbf{Solar ARs (pooled, 5)} & \textbf{magnetic} & \textbf{0.453} & \textbf{$\sim$46\%} \\
\quad AR\,11515 (most repulsive) & magnetic & 0.516 & $\sim$90\% \\
\quad AR\,12673 (giant-flare) & magnetic & 0.419 & $\sim$23\% \\
\quad AR\,11429 (least) & magnetic & 0.401 & $\sim$10\% \\
Poisson floor & --- & 0.386 & 0\% \\
\bottomrule
\end{tabular}
\end{table}

\begin{figure}[H]
\centering
\includegraphics[width=\textwidth]{figs/fig2_depletion_probe.pdf}
\caption{The reservoir-depletion probe. Single-source systems ordered by $\langle r\rangle$, which maps linearly to the reservoir depletion ratio (top axis). Earth systems (high depletion) populate the repulsion zone; solar active regions (this work, $\star$) sit at the low-depletion floor near Poisson. The failure of solar flares to repel measures where the charge-and-release mechanism switches off.}
\label{fig:probe}
\end{figure}

\subsection{Inter-AR scatter is not driven by release strength}

The five-region table shows real region-to-region scatter in $\langle r\rangle$ (0.40--0.52). A natural hypothesis, internal to the depletion picture, is that this scatter is itself a depletion signal: regions that release a larger fraction of their free energy per flare should sit higher (toward GOE). We tested this directly by extracting $\langle r\rangle$ for all 16 active regions of Solar Cycle~24 that produced $\geq6$ tagged M/X flares (enough for $\geq5$ waiting times), and correlating it against five proxies for per-event release strength: flare production rate, X-class fraction, median peak flux, maximum peak flux, and mean time-integrated flare energy.

\emph{No proxy correlates significantly with $\langle r\rangle$} (Figure~\ref{fig:robust}). The strongest association is X-class fraction (Pearson $r=+0.25$, in the predicted direction but $p=0.35$); all others are weaker, and several have the opposite sign (e.g.\ median peak flux $r=-0.29$, $p=0.27$). Every test gives $p>0.27$; none survives even a lenient threshold. The reason is structural: the per-AR bootstrap uncertainty on $\langle r\rangle$ is $\sigma\approx0.10$--0.17, comparable to the \emph{entire} Poisson--GOE gap of 0.145. At single-AR sample sizes ($n=6$--31), $\langle r\rangle$ is noise-dominated, and any underlying depletion signal is buried. This is why the apparent repulsion of an individual region such as AR\,11515 must be read as a statistical fluctuation, not a measurement. The depletion interpretation operates \emph{across} systems---where each system contributes a well-sampled pooled $\langle r\rangle$---not \emph{within} the solar AR population, where individual estimates are too noisy to order.

\begin{figure}[H]
\centering
\includegraphics[width=\textwidth]{figs/fig4_proxy_robustness.pdf}
\caption{Robustness test. Per-AR $\langle r\rangle$ (16 active regions, Solar Cycle~24) against four proxies for single-event release strength. None correlates significantly (all $p>0.27$); the weighted regression lines are flat. Point size and color encode the number of M/X flares; error bars are bootstrap $1\sigma$. Per-AR $\langle r\rangle$ is noise-dominated, so only pooled estimates are interpretable.}
\label{fig:robust}
\end{figure}

\subsection{Single-AR Poisson behavior is invariant across the solar cycle}

If single-AR Poisson behavior were an artifact of one epoch, it should vary with cycle phase. It does not. Pooling single-AR normalized waiting times within each of our four years---spanning the rise/early-maximum (2012), maximum (2014), early decline (2015), and late decline (2017) of Solar Cycle~24---gives single-AR $\langle r\rangle$ of 0.463, 0.415, 0.415, and 0.419 respectively (Figure~\ref{fig:phase}a). A weighted consistency test cannot distinguish these from a single common value ($\chi^2=0.96$, 3 d.o.f., $p=0.81$; weighted mean $\langle r\rangle=0.428$). The Poissonian character of a single active region is therefore a \emph{universal} property across the cycle, not a coincidence of any particular phase---reinforcing that low per-event depletion is intrinsic to magnetic reconnection, independent of how active the Sun is overall.

The whole-disk signal, by contrast, \emph{does} evolve with phase, and in an informative direction. The whole-disk waiting-time CV rises monotonically from 2.23 (rise) through 2.32 (maximum) and 2.54 (early decline) to 3.10 (late decline) (Figure~\ref{fig:phase}b): disk-integrated flaring becomes \emph{more} clustered as the cycle wanes. This is consistent with the depletion picture extended to the whole disk---in decline, fewer active regions are present, so the disk signal is dominated by a small number of regions producing homologous flare trains, sharpening the short-waiting-time excess. Clustering grows precisely as the multi-source averaging that would otherwise wash it out weakens.

\begin{figure}[H]
\centering
\includegraphics[width=\textwidth]{figs/fig3_cycle_phase.pdf}
\caption{Cycle-phase analysis. (a) Single-AR pooled $\langle r\rangle$ across four phases of Solar Cycle~24 is statistically constant ($\chi^2$ test $p=0.81$), always near Poisson, never approaching GOE. (b) Whole-disk waiting-time CV rises from 2.2 to 3.1 from maximum into decline: disk-integrated flaring grows more clustered as active regions become fewer. Error bars are bootstrap $1\sigma$.}
\label{fig:phase}
\end{figure}

\section{Discussion}

\subsection{A negative result that strengthens the program}

It would have been easy, and wrong, to force the solar data into the GOE narrative. The data say Poisson, and we report Poisson. Far from undermining the unifying principle, this sharpens it. The original claim---``single long-memory charge-and-release sources repel''---contained a hidden assumption: that the release substantially depletes the source. The solar test isolates and falsifies that assumption in the one regime where it fails, and in doing so reveals the true control variable. Repulsion is governed not by single-sourceness alone but by the product of single-sourceness \emph{and} high per-event depletion. The Earth systems satisfy both; solar ARs satisfy only the first.

\subsection{$\langle r\rangle$ as a remote depletion gauge}

The practical payoff is that $\langle r\rangle$ becomes a remote probe of an otherwise hard-to-measure quantity. For a system whose internal energetics are inaccessible, the waiting-time spacing ratio reports how completely each event drains the source. A system observed to have $\langle r\rangle\approx0.40$ is telling us its events barely dent its reservoir; one at $\langle r\rangle\approx0.70$ is telling us its events nearly empty it. This is a genuinely new use of the RMT diagnostic: not classification, but quantification of release physics. For solar physics specifically, it offers an independent, model-free confirmation that flares are low-depletion events---consistent with the self-organized-criticality picture in which an AR hovers near a critical state and sheds small fractions of its energy in avalanches of all sizes \citep{lu1991,aschwanden2011}.

\subsection{Two thermodynamic states: characteristic-scale oscillator vs.\ scale-free criticality}

The depletion axis maps onto a deeper distinction in statistical physics. GOE-class repulsion and self-organized criticality (SOC) are not merely two outcomes; they are two qualitatively different thermodynamic states of a driven system. A high-depletion system possesses a \emph{characteristic scale}: the recharge time sets a preferred recurrence interval, like a clock that must be fully wound before it can run again. Its event-size and waiting-time distributions are peaked, and its spacings repel. A low-depletion system is \emph{scale-free}: it sits perpetually near a critical threshold, where any small perturbation can trigger an avalanche of any size, and the system never resets. Its event sizes follow a power law (the hallmark of SOC, and precisely what is observed for solar flare energies; \citealp{lu1991,aschwanden2011}) and its waiting times carry no characteristic scale---hence Poisson, or clustered. In these terms, \emph{the transition from GOE repulsion to Poisson clustering marks the boundary between a characteristic-scale relaxation oscillator and a scale-free self-organized critical state.} The Earth systems in this program are relaxation oscillators; the solar corona is a self-organized critical system; and $\langle r\rangle$ is the order parameter that distinguishes them. The reservoir depletion ratio is the microscopic control variable, and the RMT universality class is the macroscopic state it selects.

\subsection{Why clustering, not just Poisson, for the whole disk}

The whole-disk $\langle r\rangle=0.354$ with $\mathrm{CV}=2.32$ is sub-Poissonian (clustered). Two effects combine: superposition of many independent AR clocks (which alone would give Poisson) and sympathetic/homologous triggering (which adds short-waiting-time excess). The result is the mirror image of repulsion: instead of forbidding short intervals, the disk \emph{favors} them. On the depletion axis, clustering corresponds to effectively negative depletion---an event makes the next event \emph{more}, not less, likely---placing the whole disk below the Poisson floor, at the opposite extreme from the geyser. Phenomenologically, this large-flare-triggers-short-interval-cascade behavior is a statistical parallel to the Omori-law aftershock cascades we documented in shallow seismotectonic rhythms \citep{chen2026fault}: in both cases a large release transiently raises the probability of closely following events, driving the spacing distribution below Poisson. We stress that the parallel is statistical and phenomenological---the underlying physics differs (magnetic reconnection and sympathetic triggering versus elastic stress transfer)---but the resulting departure from Poisson toward clustering is mathematically of the same character.

\subsection{Limitations}

\emph{Sample size and inter-AR scatter.} Individual-AR samples are modest ($n=11$--41), inherent to using single regions, so individual CIs are wide and some regions (AR\,11515) reach weak GOE. We address this by pooling ($n=112$), which excludes GOE at 95\% confidence. We initially hypothesized the inter-AR scatter was a depletion signal, but the formal test (Section~4.3, 16 ARs) finds no proxy for release strength predicts $\langle r\rangle$ ($p>0.27$ for all), because $\sigma_{\langle r\rangle}\approx0.10$--0.17 at single-AR sizes swamps the 0.145-wide Poisson--GOE gap. We therefore report the scatter as noise-dominated and confine the depletion interpretation to the cross-system level, where pooled estimates are robust.

\emph{Window vs tag.} AR\,12673's selection uses a time window justified by its being the sole tagged region that month; AR\,12192 uses direct AR tags. Both selections give Poisson, and the two methods agree, so the conclusion is robust to selection method.

\emph{Class threshold.} We analyze M/X flares. Including C-class raises $n$ but mixes in more background and more clustering; the whole-disk C/M/X sample (not shown in the main table) gives $\langle r\rangle=0.41$, still non-repulsive. The single-AR Poisson result is stable to reasonable threshold choices.

\emph{Non-stationarity.} An AR's flare productivity rises and falls over its disk transit, making the underlying rate non-stationary. This does not create repulsion; if anything it adds variance. A non-stationary Poisson process remains the appropriate null, and the data are consistent with it.

\subsection{The shape of the unified picture}

With solar flares added, the program now spans ten systems and a vast range of physics, and the organizing variable is clear. On one axis sits the number of sources (single vs.\ many); on the other, the per-event depletion ratio. Repulsion lives in one corner only: single-source \emph{and} high-depletion. Move to many sources and you get Poisson (or clustering if they interact); move to low depletion and you get Poisson even with a single source. The Earth systems occupy the repulsion corner; the Sun, for all its violence, does not---and that is precisely what makes it the perfect calibration point for the lower limit of the law.

\section{Conclusions}

\begin{enumerate}
\item Using real GOES flare data from the five strongest active regions of Solar Cycle~24, the pooled flare waiting-time statistics are \emph{Poissonian}, not GOE: $\langle r\rangle=0.453$ ($n=112$), with the 95\% CI excluding the GOE value. Individual regions scatter across the Poisson--GOE boundary, reflecting variable per-event depletion.
\item The whole-disk solar-maximum sample clusters below Poisson ($\langle r\rangle=0.35$, $\mathrm{CV}=2.32$), driven by sympathetic/homologous flaring.
\item This negative result identifies the missing control variable in the unifying principle: level repulsion requires not only a single source but \emph{high per-event reservoir depletion}. Solar flares fail the GOE test because a flare empties only $\lesssim$10\% of an active region's magnetic free energy.
\item The spacing ratio $\langle r\rangle$ is thereby promoted from a single-vs-multi-source classifier to a quantitative \emph{reservoir-depletion probe}: geysers ($\sim$100\%), faults/plumes/ores ($\sim$60--90\%), glaciers and the biosphere ($\sim$55--75\%), and solar active regions ($\sim$10--90\% across regions, pooled $\sim$46\%) form a monotonic depletion ladder, sitting below all Earth systems.
\end{enumerate}

The Sun does not break the law of level repulsion; it measures where the law turns off. By flushing almost nothing per flare, a solar active region marks the zero-depletion floor against which every fully flushing Earth system is calibrated.

\section*{Code and Data Availability}
All code, the extracted GOES flare sequences, and figures: \url{https://github.com/Ruqing1963/solar-flare-rmt-depletion}. Raw data: NOAA/NCEI GOES-R XRS Flare Report \citep{nceiflrpt}.\\
Companion studies: mantle plumes (\url{https://zenodo.org/records/20768420}); metallogeny (\url{https://zenodo.org/records/20768849}); cyclostratigraphy (\url{https://zenodo.org/records/20774581}); mass extinctions and radiations (\url{https://zenodo.org/records/20783763}).

\begin{thebibliography}{40}
\bibitem[Aschwanden(2011)]{aschwanden2011} Aschwanden, M.J. (2011). \textit{Self-Organized Criticality in Astrophysics}. Springer-Praxis.
\bibitem[Aschwanden et al.(2014)]{aschwanden2014} Aschwanden, M.J., et al. (2014). 25 years of self-organized criticality: solar and astrophysics. \textit{Space Sci. Rev.}, 198, 47--166.
\bibitem[At\'as et al.(2013)]{atas2013} At\'as, Y.Y., Bogomolny, E., Giraud, O., Roux, G. (2013). Distribution of the ratio of consecutive level spacings. \textit{Phys. Rev. Lett.}, 110, 084101.
\bibitem[Chen(2026a)]{chen2026strat} Chen, R. (2026a). Spatial level repulsion in the stratigraphic record: a random matrix theory test of cyclostratigraphic bed spacing. Zenodo, \url{https://zenodo.org/records/20774581}.
\bibitem[Chen(2026b)]{chen2026plume} Chen, R. (2026b). Level repulsion in the deep Earth: a random matrix theory analysis of mantle-plume eruption rhythms. Zenodo, \url{https://zenodo.org/records/20768420}.
\bibitem[Chen(2026c)]{chen2026metal} Chen, R. (2026c). Charge-and-release rhythms in ore-forming systems: a random matrix theory test of metallogenic timing across four provinces. Zenodo, \url{https://zenodo.org/records/20768849}.
\bibitem[Chen(2026d)]{chen2026evo} Chen, R. (2026d). The rhythm of life and death: random matrix theory reveals level repulsion in Phanerozoic mass extinctions and radiations. Zenodo, \url{https://zenodo.org/records/20783763}.
\bibitem[Chen(2026e)]{chen2026hydro} Chen, R. (2026e). From geysers to megafloods: random matrix theory reveals universal level repulsion in hydrogeological charge-and-release systems. GitHub, \url{https://github.com/Ruqing1963/hydro-rmt-geysers-floods}.
\bibitem[Chen(2026f)]{chen2026fault} Chen, R. (2026f). From hypothesis to reality: scale-dependent universality classes in Earth's seismotectonic rhythms. Zenodo, \url{https://zenodo.org/records/20768130}.
\bibitem[Emslie et al.(2012)]{emslie2012} Emslie, A.G., et al. (2012). Global energetics of thirty-eight large solar eruptive events. \textit{Astrophys. J.}, 759, 71.
\bibitem[Lu \& Hamilton(1991)]{lu1991} Lu, E.T., Hamilton, R.J. (1991). Avalanches and the distribution of solar flares. \textit{Astrophys. J.}, 380, L89--L92.
\bibitem[NCEI(2025)]{nceiflrpt} NOAA National Centers for Environmental Information (2025). GOES-R Series XRS L2 Flare Report (science-quality), \url{https://data.ngdc.noaa.gov/platforms/solar-space-observing-satellites/goes/}.
\bibitem[Priest \& Forbes(2002)]{priest2002} Priest, E.R., Forbes, T.G. (2002). The magnetic nature of solar flares. \textit{Astron. Astrophys. Rev.}, 10, 313--377.
\bibitem[Shibata \& Magara(2011)]{shibata2011} Shibata, K., Magara, T. (2011). Solar flares: magnetohydrodynamic processes. \textit{Living Rev. Solar Phys.}, 8, 6.
\bibitem[Wheatland(2000)]{wheatland2000} Wheatland, M.S. (2000). The origin of the solar flare waiting-time distribution. \textit{Astrophys. J.}, 536, L109--L112.
\bibitem[Wheatland(2002)]{wheatland2002} Wheatland, M.S. (2002). Flare frequency--size distributions and waiting-time statistics. \textit{Solar Phys.}, 211, 255--274.
\bibitem[Yang et al.(2017)]{yang2017} Yang, S., et al. (2017). Eruption of a multi-flux-rope system in solar active region 12673 leading to the two largest flares in Solar Cycle 24. \textit{Astrophys. J. Lett.}, 849, L21.
\end{thebibliography}

\end{document}
